% Round 16. All-degree two-component scalar correction and Gysin-image calculation.
% Root and global_height full independent reviews PASS; higher_trace separately
% PASS on the odd-central Gysin-image calculation. The final mathematical source
% after display/generic-field clarification has SHA
% b6a7b6d450732297f19023b256ef5769fb764a763203d01f958e29f664d9401f.
% This activation changes only the status notice.
% Preserve the historical higher-kernel provenance.
\section{An all-degree scalar correction and hypersurfaces in every dimension}
\label{r16sc:section}

\paragraph{Provenance and the additional calculation.}
Historical Round 7 already recorded a numerical-admissibility conclusion for
even-dimensional generic fibers with complete-intersection double
strata, even with many components:
\path{history_01_07/bb_round7/proofs/threefold_scalar_degree_kernel.tex},
\texttt{bb7tf:odd-CI-double-all-codimensions}. Its historical
review flag is retained. We give a shorter independent correction
for two components and use the new all-degree native actual descent.
For the explicit split hypersurface degeneration we also compute
the Gysin \emph{image} in the remaining odd-dimensional central
degree. That calculation closes the central case without asserting
that the full middle cohomology of the double stratum is one-dimensional.
No novelty claim relative to the literature is made.

\begin{definition}[A scalar boundary image in degree $r$]
\label{r16sc:definition}
Let $\mathcal V/R$ be an integral regular projective strictly
semistable model of relative dimension $n\geq3$ over a $p$-adic
valuation ring, with smooth geometrically integral generic fiber,
and with reduced special fiber exactly $D_1\cup D_2$.
Assume $D_1,D_2$ and $E=D_1\cap D_2\ne\varnothing$ are smooth
projective geometrically integral over the finite residue field $k$.
Fix one relatively ample line bundle $H$, put $H_i=H|_{D_i}$,
$h=H|_E$, and fix $\ell\ne p$.
For $1\leq r\leq n$, the scalar boundary-image condition means
\[
 \operatorname{im}\left(
 H^{2r-2}(E_{\bar k},\mathbf Q_\ell(r-1))
   \xrightarrow{j_{i*}}
 H^{2r}((D_i)_{\bar k},\mathbf Q_\ell(r))\right)
       =\mathbf Q_\ell\,j_{i*}c_{1,\ell}(h)^{r-1}
       \qquad(i=1,2).
 \tag{SC16.1}\label{r16sc:image}
\]
Here $j_i:E\hookrightarrow D_i$. Each displayed generator is
nonzero, since its pairing with $c_{1,\ell}(H_i)^{n-r}$ is
the positive integer
\[
                           e=\deg_k(h^{n-1})>0.
 \tag{SC16.2}\label{r16sc:degree}
\]
Full generation $H^{2r-2}(E_{\bar k},\mathbf Q_\ell(r-1))
=\mathbf Q_\ell c_{1,\ell}(h)^{r-1}$ is sufficient, but is
not required by \eqref{r16sc:image}.
\end{definition}

\begin{theorem}[Direct rational vertical correction in every degree]
\label{r16sc:correction}
Assume \eqref{r16sc:image} in degree $r$. Every rational
codimension-$r$ class $z$ on the generic fiber with zero geometric
$\ell$-adic class has a rational Chow extension whose restrictions
to both components are geometrically homologically trivial,
and therefore numerically trivial. More precisely, for any
rational Chow extension $\Theta$ put $\theta_i=i_i^!\Theta$,
where $i_i:D_i\hookrightarrow\mathcal V$. Then
\[
 c=\frac{\deg_k(\theta_1 H_1^{n-r})}{e}\in\mathbf Q,
 \qquad
       \Theta^{\rm ad}=\Theta-i_{2*}(cH_2^{r-1})
 \tag{SC16.3}\label{r16sc:explicit-correction}
\]
has these zero component classes and unchanged generic Chow class.
No higher-codimension primitive normalization is needed.
\end{theorem}

\begin{proof}
Choose the actual Chow extension first, for instance by closure of
a cycle representative. Its frozen pure operational $K_0$
representative \texttt{r12n:K0} gives a single global $\ell$-adic
degree-$2r$ Chern class on the model with component restrictions
$a_i=\operatorname{cl}_\ell(\theta_i)$ and zero geometric generic
restriction. The proper nearby-cycle comparison and subsequent
restriction to the smooth open $D_i\setminus E$, followed by
purity for the one smooth divisor, are exactly the argument of
Lemma~\ref{r16tb:numerical}, now in degree $2r$.
They show that $a_i$ is in the Gysin image in
\eqref{r16sc:image}. Thus
\[
                   a_i=c_i j_{i*}c_{1,\ell}(h)^{r-1}
                   \qquad(c_i\in\mathbf Q_\ell).
 \tag{SC16.4}\label{r16sc:scalar-classes}
\]
The projection formula and \eqref{r16sc:degree} determine
these coefficients:
\[
                  c_i=\frac{\deg_k(\theta_i H_i^{n-r})}{e}
                           \in\mathbf Q.
 \tag{SC16.5}\label{r16sc:rational-coefficients}
\]
In particular no rational algebraic lifting of an unspecified
$\ell$-adic preimage is being assumed.

Apply frozen \texttt{r12n:weighted-trace}, equation \texttt{NL16}
in \path{round12/numerical_lift/curve_lifts_and_actual_green.tex},
to the global top-degree class
$\operatorname{ch}_{r,\ell}(\Theta)\,c_{1,\ell}(H)^{n-r}$,
where the notation denotes the chosen pure operational
representative. The two special multiplicities are one.
Its geometric generic restriction is zero, so
\[
 \deg_k(\theta_1H_1^{n-r})+
 \deg_k(\theta_2H_2^{n-r})=0,
             \qquad c_1+c_2=0.
 \tag{SC16.6}\label{r16sc:sum-zero}
\]
All degrees in this equation are native degrees, not normalized
geometric sheet counts.

Because $D_1+D_2$ is the principal divisor of a uniformizer,
$N_{D_2/\mathcal V}=\mathcal O_{D_2}(-E)$. Transverse refined
restriction and the self-intersection formula give
\[
 \begin{split}
 i_1^!i_{2*}(cH_2^{r-1})&=c\,j_{1*}h^{r-1},\\
 i_2^!i_{2*}(cH_2^{r-1})&=-c\,j_{2*}h^{r-1}.
 \end{split}
 \tag{SC16.7}\label{r16sc:correction-signs}
\]
Equations \eqref{r16sc:scalar-classes}--\eqref{r16sc:sum-zero}
therefore make the two component cycle classes of
\eqref{r16sc:explicit-correction} zero. These are classes of
specified Chow restrictions. Their homological vanishing implies
numerical vanishing by the intersection trace formula; no cycle
class of an abstract numerical quotient has been introduced.
The correction is vertical, so the generic Chow class stays $z$.
\end{proof}

\begin{corollary}[Native actual activation of the scalar image condition]
\label{r16sc:actual}
Let $X/K$ be smooth projective geometrically integral. Suppose
that, after a finite extension $F/K$, a smooth projective
geometrically integral variety $V/F$ has a model satisfying
\eqref{r16sc:image} in degree $r$, and has a projective dominant
morphism $V\to X_F$. For every
$z\in\operatorname{CH}^r(X)_{\mathbf Q,\mathrm{geom\,hom}}$,
\[
                         \mathfrak o_{X,K}^r(z)=0.
 \tag{SC16.8}\label{r16sc:actual-zero}
\]
Every specified cycle in this class has a strong Green form on
its exact support complement, with every original compact actual
test. If the source model satisfies the image condition in all
degrees, the target has full strong G in every positive codimension.
\end{corollary}

\begin{proof}
Apply Theorem~\ref{r16sc:correction} to the refined generic
pullback class. Frozen \texttt{r12n:K0} applied to its corrected
Chow lift gives the rational virtual-bundle witness whose degree-$r$
character has zero supported degree on every vertical integral
$r$-fold. Corollary~\ref{r16G:numerical-witness} applies in the
original $K$-actual theory to $V\to X_F\to X$. It retains the
ramification and residue factors of the model viewed over
$K^\circ$ and the complete degrees in dominant Bott--Chern descent.
The exact-support construction of
Proposition~\ref{r16G:dominant-witness} then gives the statement
for the originally specified cycle and the full original test space.
The source and target need not have equal dimension. No rational
map, arbitrary proper actual pushforward, or finite-on-complement
property is being assumed.
\end{proof}

\begin{lemma}[The codimension-two etale blowup decomposition]
\label{r16sc:etale-blowup}
Let $D$ be smooth projective over an algebraically closed field,
let $S\subset D$ be smooth of codimension two, let
$b:D'\to D$ be its blowup, and let $j:F\hookrightarrow D'$,
$p:F\to S$ be the exceptional divisor and its projection.
For $\ell$ invertible in the field, and with all usual Tate
twists retained, the map
\[
 b^*+j_*p^*:
 H^q(D,\mathbf Q_\ell(t))\oplus
 H^{q-2}(S,\mathbf Q_\ell(t-1))
               \longrightarrow H^q(D',\mathbf Q_\ell(t))
 \tag{SC16.9}\label{r16sc:blowup-cohomology}
\]
is an isomorphism.
\end{lemma}

\begin{proof}
For completeness the exact maps needed here are verified from
the primary blowup-square triangle: Stacks Project,
Section 59.105, Lemmas 59.105.1--59.105.2, Tag \texttt{0EW4},
\url{https://stacks.math.columbia.edu/tag/0EW4}.
It gives the long exact sequence with successive terms
\[
 H^q(D)\longrightarrow H^q(D')\oplus H^q(S)
       \xrightarrow{\,j^*-p^*\,}H^q(F)
                   \longrightarrow H^{q+1}(D).
\]
These statements may be applied to finite $\ell$-power coefficients
and then passed to rational $\ell$-adic cohomology; the proper
smooth groups are finite and give the usual Mittag--Leffler passage.
Write $\xi=c_1(\mathcal O_F(1))$ in the lines convention.
The projective-bundle formula and divisor self-intersection give
\[
 H^q(F)=p^*H^q(S)\oplus\xi p^*H^{q-2}(S)(-1),
                    \qquad j^*j_*p^*v=-\xi p^*v.
\]
Thus $j^*-p^*$ is surjective in every degree: its first
summand is supplied by $S$ and its second by $j_*p^*$.
The long exact sequence is consequently short exact in each degree.
If $j^*\omega=p^*u+\xi p^*v$, then
$j^*(\omega+j_*p^*v)=p^*u$. Exactness writes this pair as
$(b^*w,w|_S)$, proving spanning by the two maps in
\eqref{r16sc:blowup-cohomology}.
If $b^*w+j_*p^*v=0$, restriction to $F$ and the projective
decomposition first give $v=0$ and $w|_S=0$. Injectivity of
the initial map into $H^q(D')\oplus H^q(S)$ then gives $w=0$.
This proves the asserted decomposition and its concrete generators.
\end{proof}

\begin{proposition}[All-degree boundary images for the split hypersurface model]
\label{r16sc:hypersurface-images}
Fix $n\geq3$. Replace $\mathbf P^5_R$ in
Proposition~\ref{r16tb:hypersurface-model} by
$\mathbf P^{n+1}_R$, with the same degrees $a,b,a+b$.
Assume that $D_1=V(\bar A)$ and $D_2=V(\bar B)$ are smooth
geometrically integral $n$-folds, that
$E=V(\bar A,\bar B)$ is a smooth geometrically integral
$(n-1)$-fold, and that
$\Sigma=V(\bar A,\bar B,\bar G)$ is a smooth geometrically
integral $(n-2)$-fold, with the indicated transverse differentials.
Assume the generic fiber
$V=V(AB+\pi G)\subset\mathbf P^{n+1}_{F}$ is smooth and
geometrically integral.
Then the model obtained by blowing up $(A,\pi)$ satisfies
the scalar boundary-image condition in \emph{every} degree
$1\leq r\leq n$, for any chosen relatively ample $H$.
\end{proposition}

\begin{proof}
The local equation and the saturated charts in
\eqref{r16tb:charts} are unchanged; there are now $n-2$
additional free ambient coordinates instead of two.
The same chart proof gives an integral regular projective
strictly semistable model with exactly two reduced components
\[
 D'_1=\operatorname{Bl}_{\Sigma}D_1,\qquad D'_2=D_2,
                       \qquad D'_1\cap D'_2=E.
 \tag{SC16.10}\label{r16sc:general-model}
\]
The exceptional locus over $\Sigma$ has dimension $n-1$
and codimension two in the total model, so it is not an
additional special-fiber component. The normal bundles along
the double stratum are still $\mathcal O_E(-a)$ and
$\mathcal O_E(a)$.

If $2r-2\ne n-1$, smooth complete-intersection cohomology
gives
\[
 H^{2r-2}(E_{\bar k},\mathbf Q_\ell(r-1))
                         =\mathbf Q_\ell h_0^{r-1},
 \tag{SC16.11}\label{r16sc:nonmiddle}
\]
where $h_0=c_{1,\ell}(\mathcal O_E(1))$.
The primary statement is Deligne, SGA 7 II, Expose XI,
Theorem 1.6(ii)--(iii), printed p.~48, checked in the
original French scan cited in Section~\ref{r16tb:section}.
Below the middle it is weak Lefschetz; above the middle
the stated theorem, or Poincare duality, gives the same
rational generator. The power of any ample restriction
$h=H|_E$ is nonzero, as its complementary ample degree
is positive, so it also generates this one-dimensional
space. This proves \eqref{r16sc:image} in these degrees.
In particular these are all degrees when $n$ is even.

It remains to treat $n=2m-1$ and $r=m$, with $m\geq2$.
The full middle cohomology $H^{n-1}(E)$ may be large.
Poincare duality identifies the rank of
\[
 j_{i*}:H^{n-1}(E)\longrightarrow H^{n+1}(D'_i)
\]
with the rank of its adjoint restriction
$H^{n-1}(D'_i)\to H^{n-1}(E)$; here and in this rank
calculation the compatible Tate twists are suppressed.
For $D'_2=D_2$, weak Lefschetz gives
$H^{n-1}(D_2)=\mathbf Q_\ell H_0^{m-1}$, and its restriction
is the nonzero class $h_0^{m-1}$. The restriction rank is one.

For $D'_1$ write $b:D'_1\to D_1$ and let
$j_F:F\hookrightarrow D'_1$, $p:F\to\Sigma$ be the
exceptional divisor. Lemma~\ref{r16sc:etale-blowup} gives
\[
 H^{n-1}(D'_1)
    =b^*H^{n-1}(D_1)
            \oplus j_{F*}p^*H^{n-3}(\Sigma)(-1).
 \tag{SC16.12}\label{r16sc:central-blowup}
\]
Both groups on the right are one-dimensional. Indeed the
dimensions of the smooth complete intersections $D_1$
and $\Sigma$ are respectively the odd numbers $n$ and
$n-2$, and the two cohomological degrees lie one below
their respective middle degrees. Thus the two displayed
generators can be taken as
\[
                b^*H_0^{m-1},\qquad
                j_{F*}p^*H_{\Sigma,0}^{m-2}.
\]
Their restrictions to $E$ are, respectively,
\[
                     h_0^{m-1},\qquad
                     (a+b)h_0^{m-1}.
 \tag{SC16.13}\label{r16sc:central-restrictions}
\]
For the second formula, the strict transform of $E$ meets
$F$ transversely in $\Sigma$. In the first local chart,
inside $D'_1$ they are $v=0$ and $g=0$, respectively.
The transverse restriction--Gysin formula therefore gives
the pushforward from $\Sigma\hookrightarrow E$ of
$H_{\Sigma,0}^{m-2}$. That divisor has class
$(a+b)h_0$, proving the formula.
Consequently the restriction image from $D'_1$ also has
dimension one, and Poincare duality proves that each of
the two relevant Gysin images has dimension one.

For the actual chosen relative ample restriction $h$, the
element $j_{i*}c_{1,\ell}(h)^{m-1}$ is nonzero by
\eqref{r16sc:degree}. It therefore generates the respective
Gysin image, proving \eqref{r16sc:image} in the central
degree as well. When $n=3$ we do not assert that the full
$H^2(E)$ is one-dimensional or that $h$ is a multiple of
$h_0$; nonvanishing in a one-dimensional \emph{image}
is the only assertion needed.
\end{proof}

\begin{theorem}[Full local G for the supplied split hypersurface families]
\label{r16sc:full-G}
For every $n\geq3$ and every pair $a,b\geq1$, each smooth
projective hypersurface supplied with the transverse local
degeneration of Proposition~\ref{r16sc:hypersurface-images}
satisfies full strong G in every positive codimension.
The same is true for every smooth projective geometrically
integral dominant image of such a hypersurface after a finite
valued extension, with all degrees, forms and actual tests
interpreted over the original local field.

For each fixed $n,a,b$, the indicated local family is nonempty
after a finite unramified extension of any given $p$-adic
field. There is no bound on $a+b$. In particular the construction
includes degrees $a+b>n+2$, where the smooth hypersurface has
ample canonical bundle.
\end{theorem}

\begin{proof}
Proposition~\ref{r16sc:hypersurface-images} supplies the
scalar image condition in every degree. Apply
Theorem~\ref{r16sc:correction} and
Corollary~\ref{r16sc:actual} degree by degree. These prove
zero obstruction, and hence exact-support strong G, not merely
a current identity or a statement for algebraic test classes.

The first-jet incidence and residue-coset density proof of
Corollary~\ref{r16tb:nonempty} works in
$\mathbf P^{n+1}$ with the same dimensions shifted.
Here $\dim\Sigma=n-2\geq1$, so the smooth complete intersections
are geometrically connected, and hence geometrically integral,
as in that proof. A finite residue extension defines the
transverse data; their lifts can be chosen with smooth generic
hypersurface because its smooth-coefficient locus is nonempty
Zariski open and the fixed residue coset of $G$ is Zariski dense.
Finally adjunction for a degree-$d$ smooth hypersurface in
$\mathbf P^{n+1}$ gives $\omega_V=\mathcal O_V(d-n-2)$,
which is ample when $d=a+b>n+2$.
\end{proof}

\begin{corollary}[Global height on the corresponding placewise class]
\label{r16sc:height}
Let $X$ be smooth projective geometrically integral over a number
field. Suppose that at every finite place its full homological
Chow groups are covered by Theorem~\ref{r16sc:full-G}, possibly
through its finite-extension dominant-witness clause, or by
another proved full obstruction-zero local criterion.
Then the canonical real height of Theorem~\ref{r16h:global-height}
is defined on the full complementary geometrically homologically
trivial rational Chow groups of $X$.
\end{corollary}

\begin{proof}
The proved local statements give zero actual curvature obstruction
for every required class at every finite place. Thus the local
obstruction-zero domain in Theorem~\ref{r16h:global-height}
is the full geometric homological Chow group in each degree.
Apply that theorem, including its archimedean construction,
native field normalization, principal law and good-place finiteness.
No converse from the existence of a Green form to obstruction
vanishing is needed.
\end{proof}

\paragraph{Limits.}
No assertion is made that an arbitrary hypersurface admits this
two-component degeneration after extension, or that an arbitrary
variety is dominated by one of these supplied models. The central
odd-dimensional argument used the particular hypersurface and
complete-intersection center in \eqref{r16sc:general-model}; it
does not close the middle-degree gate for arbitrary two-component
SNC models. The global consequence requires its stated local
hypotheses at every finite place and is not a proof of the general
unconditional Beilinson--Bloch height conjecture.
